\documentclass[11pt]{article}
\usepackage[a4paper,margin=25mm]{geometry}
\usepackage[T1]{fontenc}
\usepackage{lmodern,amsmath,amssymb,amsthm,mathtools}
\usepackage[colorlinks=true,linkcolor=blue,urlcolor=blue,citecolor=blue]{hyperref}
\newtheorem{theorem}{Theorem}
\newtheorem{lemma}[theorem]{Lemma}
\newcommand{\had}{\operatorname{had}}
\newcommand{\PP}{\mathcal P}
\title{A component transfer theorem for Hadwiger--Ramsey numbers}
\author{}
\date{}
\begin{document}
\maketitle

\section{The minor bounds and component optimization}

Write $\had(G)$ for the largest integer $t$ such that $G$ contains a $K_t$
minor, and $\kappa(G)$ for its vertex connectivity. All logarithms are natural.
We use the following two bounds of Norin and Rahman.

\begin{theorem}[Theorem 1.4 from {\cite{norinrahman}}]\label{thm:random}
For every $\varepsilon>0$ there exists $n_0=n_0(\varepsilon)$ such that the
following holds for every $n\geq n_0$. Let $(X_1,\ldots,X_k)$ be a partition
of $[n]$, and let $0\leq p_{ij}=p_{ji}\leq1-\varepsilon$ for all
$i,j\in[k]$. Form a random graph $G$ on $[n]$ by including each edge $uv$,
with $u\in X_i$ and $v\in X_j$, independently with probability $p_{ij}$.
Put $\alpha_i=|X_i|/n$ and
\[
 L=\sqrt{-\sum_{i,j=1}^k\alpha_i\alpha_j\log(1-p_{ij})}.
\]
If $L>\varepsilon$, then
\[
 \mathbb P\left(\had(G)\leq(1+\varepsilon)L\frac{n}{\sqrt{\log n}}\right)
 \geq1-\varepsilon.
\]
\end{theorem}

\begin{theorem}[Theorem 1.5 from {\cite{norinrahman}}]\label{thm:minor}
For every $\varepsilon>0$ there exists $n_0=n_0(\varepsilon)$ such that the
following holds for every $n\geq n_0$. Let $G$ be an $n$-vertex graph with
$\kappa(G)\geq\varepsilon n$, and let $(X_1,\ldots,X_k)$ be a partition of
$V(G)$. Put
\[
 \alpha_i=\frac{|X_i|}{n},\qquad
 p_{ij}=\frac{e_G(X_i,X_j)}{|X_i||X_j|}\quad(i\ne j),\qquad
 p_{ii}=\frac{e(G[X_i])}{\binom{|X_i|}{2}},
\]
and $p_{ij}^*=\min\{p_{ij},1-\varepsilon\}$.
If $\alpha_i\geq\varepsilon$ for every $i\in[k]$, then
\[
 \had(G)\geq\frac{n}{\sqrt{\log n}}
 \left(\sqrt{-\sum_{i,j=1}^k\alpha_i\alpha_j
                    \log(1-p_{ij}^*)}-\varepsilon\right).
\]
\end{theorem}

Fix an integer $r\geq2$. A \emph{component instance} consists of an integer
$k\geq1$, positive masses $\alpha_1,\ldots,\alpha_k$ summing to one,
partitions $\PP_1,\ldots,\PP_r$ of $[k]$, and symmetric matrices
$(p^c_{ij})_{i,j\in[k]}$ with entries in $[0,1)$ such that
\[
 \sum_{c=1}^r p^c_{ij}=1,
 \qquad
 p^c_{ij}=0\quad\text{if $i,j$ belong to different parts of $\PP_c$}.
\]
Define
\begin{equation}\label{eq:program}
 \beta_r=\inf\;\max_{\substack{c\in[r]\\ C\in\PP_c}}
       \sqrt{-\sum_{i,j\in C}\alpha_i\alpha_j\log(1-p^c_{ij})},
\end{equation}
where the infimum is over all component instances. All double sums are over
ordered pairs. The one-part instance with $p^c_{11}=1/r$ shows that $\beta_r$ is
finite. Also $\beta_r\geq1/r>0$. Indeed, for an instance with objective
value $M$, put
\[
 s_C=\sum_{i\in C}\alpha_i,\qquad
 d_{c,C}=\sum_{i,j\in C}\alpha_i\alpha_jp^c_{ij}.
\]
Since $-\log(1-p)\geq p$, we have
$d_{c,C}\leq\min\{M^2,s_C^2\}\leq Ms_C$. Summing over all colours and
their parts gives $1\leq rM$, proving the claimed lower bound.

\begin{theorem}\label{thm:transfer}
For every integer $r\geq2$ and every $\varepsilon>0$, all sufficiently large
$r$-edge-coloured complete graphs $K_n$ have a colour class $G_c$, $c\in[r]$, satisfying
\[
 \had(G_c)\geq(\beta_r-\varepsilon)\frac{n}{\sqrt{\log n}}.
\]
\end{theorem}

\begin{theorem}\label{thm:construction}
For every integer $r\geq2$ and every $\varepsilon>0$, all sufficiently large
$n$ admit an $r$-edge-colouring of $K_n$ whose colour classes $G_1,\ldots,G_r$
satisfy
\[
 \had(G_c)\leq(\beta_r+\varepsilon)\frac{n}{\sqrt{\log n}}
 \qquad\text{for every }c\in[r].
\]
\end{theorem}

Let $R_h(k;r)$ be the least integer $n$ such that every $r$-edge-colouring
of $K_n$ contains a monochromatic $K_k$ minor.
Theorems~\ref{thm:transfer} and~\ref{thm:construction} give the following.

\begin{theorem}\label{thm:ramsey}
For every fixed integer $r\geq2$,
\[
 R_h(k;r)=\left(\beta_r^{-1}+o(1)\right)k\sqrt{\log k}.
\]
\end{theorem}

\section{The random construction}

\begin{proof}[Proof of Theorem~\ref{thm:construction}]
Fix a component instance with objective value $M<\beta_r+\varepsilon/2$.
For each colour $c$ and each $C\in\PP_c$, put
\[
 s_C=\sum_{i\in C}\alpha_i,\qquad
 E_{c,C}=\sqrt{-\sum_{i,j\in C}\alpha_i\alpha_j\log(1-p^c_{ij})}.
\]
Thus $s_C>0$ and $E_{c,C}\leq M$. Let $q=\sum_{c=1}^r|\PP_c|$.
Since the instance is finite and all its densities are strictly less than
one, we may choose $\delta>0$ such that
\[
 q\delta<1,\qquad
 p^c_{ij}\leq1-\delta\quad\text{for all }c,i,j,\qquad
 (1+\delta)M<\beta_r+\varepsilon,
\]
and $2\delta<E_{c,C}/s_C$ for every pair $(c,C)$ with $E_{c,C}>0$.

For every sufficiently large $n$, partition $[n]$ into $k$ sets
$X_1,\ldots,X_k$ of sizes $n_i=\alpha_in+O(1)$. Independently for each
edge $uv$ with $u\in X_i$ and $v\in X_j$, assign colour $c$ with probability
$p^c_{ij}$. These probabilities sum to one. For a fixed colour $c$, its edge
indicators are independent, and there are no colour-$c$ edges between
$U_C=\bigcup_{i\in C}X_i$ and $U_D$ for distinct $C,D\in\PP_c$.

Write $N_C=|U_C|$ and
\[
 E_{c,C,n}=\sqrt{-\sum_{i,j\in C}
              \frac{n_in_j}{n^2}\log(1-p^c_{ij})},\qquad
 L_{c,C,n}=\frac{n}{N_C}E_{c,C,n}.
\]
Then $N_C/n\to s_C$, $E_{c,C,n}\to E_{c,C}$, and
$L_{c,C,n}\to E_{c,C}/s_C$. If $E_{c,C}>0$, then
$L_{c,C,n}>\delta$ and $N_C\geq n_0(\delta)$ for all sufficiently large $n$.
Theorem~\ref{thm:random}, applied to $G_c[U_C]$ with its partition
$(X_i)_{i\in C}$ and parameter $\delta$, gives
\[
 \had(G_c[U_C])\leq(1+\delta)L_{c,C,n}
                         \frac{N_C}{\sqrt{\log N_C}}
\]
with probability at least $1-\delta$. On this event,
\[
 \frac{\had(G_c[U_C])\sqrt{\log n}}{n}
 \leq(1+\delta)E_{c,C,n}\sqrt{\frac{\log n}{\log N_C}}
 <\beta_r+\varepsilon
\]
for all sufficiently large $n$: the middle expression tends to
$(1+\delta)E_{c,C}\leq(1+\delta)M<\beta_r+\varepsilon$.

If $E_{c,C}=0$, positivity of the masses forces $p^c_{ij}=0$ for all
$i,j\in C$. In that case $G_c[U_C]$ is edgeless and has Hadwiger number
$1$, which also satisfies the required bound for sufficiently large $n$.
A union bound over the at most $q$ positive-energy pairs gives simultaneous
success with probability at least $1-q\delta>0$. Choose a colouring in this
event. Since the sets $U_C$, $C\in\PP_c$, have no colour-$c$ edges between
them,
\[
 \had(G_c)=\max_{C\in\PP_c}\had(G_c[U_C])
 \leq(\beta_r+\varepsilon)\frac{n}{\sqrt{\log n}}
 \qquad(c\in[r]),
\]
as required.
\end{proof}

\section{Connectivity and graph limits}

\begin{lemma}\label{lem:core}
Fix $0<\rho<1/4$ and $\sigma>0$. Suppose an $N$-vertex graph $H$ satisfies
\[
 e_H(A,V(H)\setminus A)\geq\sigma N^2
 \quad\text{whenever}\quad
 \rho N\leq|A|\leq(1-\rho)N.
\]
For sufficiently large $N$, there is an induced subgraph of $H$ on more than
$(1-\rho)N$ vertices with vertex connectivity at least
$\lfloor\sigma N/2\rfloor$.
\end{lemma}
\begin{proof}
Put $t=\lfloor\sigma N/2\rfloor$. A cut with one side of size
$\lfloor N/2\rfloor$ gives $\sigma\leq1/4$, so $t\leq N/8$.
Starting with $H$, whenever the current induced graph is not $t$-connected,
take a separation with nonempty sides $A,B$, $|A|\leq|B|$, and separator
$Z$ of size less than $t$. Delete $A$ and retain $Z\cup B$. Fewer than
$t|A|$ edges join the deleted vertices to those retained.

Suppose the set $D$ of deleted vertices first reaches size $\rho N$.
Before this step fewer than $\rho N$ vertices had been deleted, and the
last deleted side has size at most $N/2$. Hence
\[
 \rho N\leq|D|<\rho N+N/2<(1-\rho)N.
\]
Every edge from $D$ to its complement was counted when its endpoint in $D$
was deleted. Thus
\[
 e_H(D,V(H)\setminus D)<t|D|\leq\frac{\sigma N^2}{2},
\]
contrary to the cut hypothesis. Therefore fewer than $\rho N$ vertices are
ever deleted. The current graph always has more than $3N/4>t$ vertices,
so every failure of $t$-connectivity admits the required separation. The
process terminates with the desired induced subgraph.
\end{proof}

An \emph{$r$-colour graphon} is a tuple $(W_1,\ldots,W_r)$ of symmetric
measurable functions $[0,1]^2\to[0,1]$ with $\sum_{c=1}^r W_c=1$ almost everywhere.
Write $\mu$ for Lebesgue measure. Coloured graph-limit compactness
\cite{decorated} gives a convergent subsequence of every sequence of
$r$-edge-coloured complete graphs of orders tending to infinity. Convergence
is in the cut metric, with a common alignment for all colours.

A graphon $W$ on a probability space is \emph{connected} if
$\int_{A\times A^{\mathrm c}}W>0$ whenever $0<\mu(A)<1$. By the component decomposition
\cite[Lemma 5.5]{janson}, each $W_c$ admits a countable measurable partition
\begin{equation}\label{eq:components}
 [0,1]=Z_c\mathbin{\dot\cup}\bigcup_{j\geq1}S_{c,j}
\end{equation}
up to null sets, where each $S_{c,j}$ has positive measure and supports a
connected restriction, and $W_c$ vanishes almost everywhere outside
$\bigcup_j S_{c,j}^2$. Restrictions are viewed with normalized measure.

If graphs $H_n$ converge to a connected graphon, then for every $\rho>0$
there is $\sigma>0$ such that, for all sufficiently large $n$, every cut of
$H_n$ whose sides have size at least $\rho|H_n|$ has at least
$\sigma|H_n|^2$ edges \cite[Theorem 1.10(i)]{janson}.

\begin{lemma}\label{lem:realize}
Suppose $r$-edge-coloured graphs of orders $n\to\infty$ converge to
$(W_1,\ldots,W_r)$. Every fixed finite measurable partition of $[0,1]$ into
sets of positive measure can be realized by vertex partitions such that
relative part sizes and edge densities between parts converge to their
graphon values. For every union of parts, the induced coloured graphs
converge to the corresponding restricted graphon with normalized measure.
\end{lemma}
\begin{proof}
Align the step graphons to the limit in cut norm. Their vertex cells have
measure $1/n$. Independently assign each vertex to a partition class, using
the proportions of its cell lying in the measurable classes as probabilities.

Fix a finite coloured graph with a specified partition class for each
vertex. Its expected normalized count under these assignments differs by
$O(1/n)$ from the corresponding integral in the aligned step graphons:
only maps with repeated vertices contribute an error. The cut-norm counting
lemma, with the class indicators as vertex weights, makes these integrals
converge to their values in the limit. Changing one vertex assignment alters
any fixed normalized count by $O(1/n)$, so bounded differences give
concentration. A diagonal choice over the countably many such test graphs
produces deterministic partitions for which all these counts converge.
Single-vertex and single-edge tests give the part sizes and densities;
summing over the permitted vertex classes and normalizing gives convergence
of every induced restriction.
\end{proof}

\section{Proof of the transfer theorem}

\begin{lemma}\label{lem:lift}
Suppose $r$-edge-coloured graphs of orders $n\to\infty$ converge to
$(W_1,\ldots,W_r)$, and let $G_{c,n}$ be their colour-$c$ graphs.
Let $S$ be a positive-measure connected component of $W_c$, and let
$A_1,\ldots,A_k\subseteq S$ be disjoint measurable sets of positive measure.
Put
\[
 \lambda_i=\mu(A_i),\qquad
 p_{ij}=\frac{1}{\lambda_i \lambda_j}\int_{A_i\times A_j}W_c.
\]
Then
\begin{equation}\label{eq:lift}
 \liminf_{n\to\infty}\frac{\had(G_{c,n})\sqrt{\log n}}{n}
 \geq\sqrt{-\sum_{i,j=1}^k \lambda_i \lambda_j\log(1-p_{ij})},
\end{equation}
where $-\log0=+\infty$.
\end{lemma}
\begin{proof}
Apply Lemma~\ref{lem:realize} to the partition consisting of the $A_i$,
$S\setminus\bigcup_{i=1}^k A_i$, and $[0,1]\setminus S$, omitting null parts.
Let $S_n$ and $A_{i,n}$ be the corresponding vertex sets. Then
$|S_n|/n\to\mu(S)>0$, and $G_{c,n}[S_n]$ converges to the connected
restriction of $W_c$ to $S$.

Let $\eta$ be the smallest measure of a positive-measure part of this
partition lying in $S$. Fix $0<\zeta<1$, $\nu>0$, and
$0<\rho<\min\{1/4,\eta/4\}$. The connected-limit cut bound and
Lemma~\ref{lem:core} give induced subgraphs $H_n\subseteq G_{c,n}[S_n]$
such that
\[
 |S_n\setminus V(H_n)|<\rho n,\qquad
 \kappa(H_n)\geq\gamma n
\]
for some fixed $\gamma>0$ and all sufficiently large $n$.
Write $N_n=|H_n|$. Intersect the realized partition of $S_n$ with $V(H_n)$.
Every resulting part has size at least $\eta n/2$, and hence relative size
at least $\eta/2$ in $H_n$.

For $i\leq k$, write $u_{i,n}=|A_{i,n}\cap V(H_n)|/n$, and let $q_{ij,n}$
be the densities between these parts in $H_n$, with diagonal densities
defined as in Theorem~\ref{thm:minor}. Then
\[
 u_{i,n}=\lambda_i+O(\rho)+o(1),\qquad
 q_{ij,n}=p_{ij}+O(\rho)+o(1).
\]
The constants depend only on the fixed sets $A_i$; diagonal densities
have an additional $o(1)$ correction from excluding loops.

Set $f_\zeta(p)=-\log(1-\min\{p,1-\zeta\})$, and choose
\[
 0<\delta\leq\min\{\gamma,\eta/2,\zeta,\nu\}.
\]
Apply Theorem~\ref{thm:minor} to $H_n$ with parameter $\delta$ and the full
intersected partition. Its connectivity and part-size hypotheses hold
because $N_n\leq n$. Since $f_\delta\geq f_\zeta$, retaining only the
terms corresponding to $A_1,\ldots,A_k$ gives
\[
 \frac{\had(G_{c,n})\sqrt{\log n}}{n}
 \geq\sqrt{\frac{\log n}{\log N_n}}
 \left(\sqrt{\sum_{i,j=1}^k u_{i,n}u_{j,n}f_\zeta(q_{ij,n})}
             -\delta\frac{N_n}{n}\right).
\]
Here the factors $N_n$ cancel against the normalization of the partition
masses. Since $N_n$ is bounded below by a positive multiple of $n$ and
$f_\zeta$ is bounded and Lipschitz, for a constant $C_\zeta$ independent
of $\rho$ and $\nu$ we obtain
\[
 \liminf_{n\to\infty}\frac{\had(G_{c,n})\sqrt{\log n}}{n}
 \geq\sqrt{\max\left\{0,
            \sum_{i,j=1}^k \lambda_i \lambda_j f_\zeta(p_{ij})-C_\zeta\rho\right\}}
       -\nu.
\]
Let $\nu\downarrow0$, then $\rho\downarrow0$, and finally
$\zeta\downarrow0$. Monotone convergence of the finite sum proves
\eqref{eq:lift}, including the case of an infinite right-hand side.
\end{proof}

\begin{proof}[Proof of Theorem~\ref{thm:transfer}]
The conclusion is immediate if $\beta_r\leq\varepsilon$. Otherwise, if it
fails, there are $r$-edge-coloured complete graphs of orders $n\to\infty$
and a constant $0\leq B<\beta_r$ such that
\[
 \max_{c\in[r]}\frac{\had(G_{c,n})\sqrt{\log n}}{n}\leq B.
\]
Pass to a convergent subsequence with limit $(W_1,\ldots,W_r)$, and use
the component partitions \eqref{eq:components}. Their common refinement
is countable; its positive-measure atoms $A_i$ cover $[0,1]$ up to a null
set. Put
\[
 \lambda_i=\mu(A_i),\qquad
 p^c_{ij}=\frac{1}{\lambda_i \lambda_j}\int_{A_i\times A_j}W_c.
\]
For each colour, give an atom the label of the part of
\eqref{eq:components} containing it, retaining $Z_c$ as one label.
The densities sum to one and vanish between different labels; every
colour-$c$ density involving an atom in $Z_c$ is zero.

All $p^c_{ij}$ are strictly less than one. Indeed, if $p^c_{ij}=1$, the
atoms concerned lie in a common connected component of $W_c$, and
Lemma~\ref{lem:lift} applied to these atoms gives an infinite lower bound,
contradicting $B<\infty$.

Choose finitely many atoms whose total measure $s$ satisfies $s \beta_r>B$.
Keep their densities and labels and normalize their masses to $\lambda_i/s$.
This is a component instance, so \eqref{eq:program} supplies a colour $c$
and one of its label classes $C$ for which
\[
 -\sum_{i,j\in C}\lambda_i \lambda_j\log(1-p^c_{ij})
 \geq s^2\beta_r^2>B^2.
\]
This class has positive energy and therefore lies in a connected component
of $W_c$, rather than in $Z_c$. Lemma~\ref{lem:lift} applied to its finitely
many atoms contradicts the upper bound $B$.
\end{proof}

\begin{thebibliography}{9}
\bibitem{norinrahman}
Sergey Norin and Neil Rahman,
\emph{Asymptotics of the three-color Hadwiger-Ramsey numbers}, manuscript.

\bibitem{decorated}
L\'aszl\'o Lov\'asz and Bal\'azs Szegedy,
\emph{Limits of compact decorated graphs}, arXiv:1010.5155, 2010.
\url{https://arxiv.org/abs/1010.5155}.

\bibitem{janson}
Svante Janson,
\emph{Connectedness in graph limits}, arXiv:0802.3795, 2008.
\url{https://arxiv.org/abs/0802.3795}.
\end{thebibliography}
\end{document}
